\documentclass[11pt]{article}
\usepackage[margin=1in]{geometry}
\usepackage{amsmath,amssymb,amsthm}
\newtheorem{theorem}{Theorem}[section]
\newtheorem{lemma}[theorem]{Lemma}
\newtheorem{remark}[theorem]{Remark}
\title{The Unrestricted Fifth-Order Deficit Above One Radian}
\author{}\date{October 5, 2026 --- paper-proof draft}
\begin{document}\maketitle
\begin{abstract}
For the prescribed-net-angle sofa problem, the exact regime $m(\omega)=1+\omega^2/2$ ends at $\omega=1$. We prove the first nonzero deficit immediately afterwards:
\[
 m(1+\delta)=1+\frac{(1+\delta)^2}{2}-\frac{229}{1920}\delta^5+o(\delta^5).
\]
The lower bound is supplied by actual feasible sofas, with a repaired endpoint derivative. The upper bound applies to arbitrary global maximizing caps, using the previously proved all-maximizer geometry and an exact convex transition-model certificate. Their rescaled endpoint support differences have a unique limiting profile. This is an asymptotic theorem, not an exact optimizer classification on the remaining positive-length angle interval, and it has not been independently refereed or formalized.
\end{abstract}

\section{A local coercivity lemma}
For $F,G\in H^1(0,w)$ with $F(w)=G(0)=0$, put
\[
 J_w(F,G)=\frac12\int_0^w(F'^2+G'^2+F G'-G F').
\]
This is the exact gap from Baek's normalized relaxed cap when $F,G$ are active support differences.

\begin{lemma}[Uniform coercivity and endpoint localization]\label{lem:local}
For $w$ in a sufficiently small fixed neighborhood of one, there is a constant $c>0$ such that
\[
 J_w(F,G)\ge c\bigl(\|F'\|_2^2+\|G'\|_2^2\bigr).
\]
For $0<2\eta<w$ one also has
\[
 J_w(F,G)\ge\frac12\int_0^\eta F'^2+\frac12\int_{w-\eta}^wG'^2
 -C\eta\bigl(\|F'\|_2^2+\|G'\|_2^2\bigr)
\]
with a uniform constant $C$.
\end{lemma}
\begin{proof}
Integration by parts and the endpoint traces give
\[
 J_w=\frac12\left(\|F'\|_2^2+\|G'+F\|_2^2-\|F\|_2^2\right).
\]
The one-ended Poincare inequality gives $\|F\|_2^2\le4w^2\|F'\|_2^2/\pi^2$. Since $w$ stays below $\pi/2$ by a fixed margin, this bounds $\|F'\|_2$ and $\|G'+F\|_2$ in terms of $J_w^{1/2}$. The triangle inequality then controls $\|G'\|_2$, proving uniform coercivity. In particular $J_w$ is nonnegative for every pair with the stated traces, whether or not the pair comes from cap supports.

For localization set $F_\ell(t)=F(t)-F(\eta)$ on $[0,\eta]$ and zero afterwards, and set $G_\ell(t)=G(t)-G(w-\eta)$ on $[w-\eta,w]$ and zero before. Put $F_b=F-F_\ell$, $G_b=G-G_\ell$. The derivative energies split orthogonally, and the two localized supports are disjoint. Thus $J_w(F_\ell,G_\ell)$ is exactly the two displayed endpoint derivative energies. The bulk pair retains the pinned traces, so its $J_w$ is nonnegative.

The remaining bilinear terms reduce by integration by parts to integrals of $F_\ell G_b'$ and $G_\ell F_b'$, with signs irrelevant for the bound. Endpoint terms vanish: $G_b(0)=F_b(w)=0$ and the localized functions vanish at their other support endpoints. The estimates $\|F_\ell\|_2\le\eta\|F'\|_2$ and $\|G_\ell\|_2\le\eta\|G'\|_2$ bound these terms by $C\eta(\|F'\|_2^2+\|G'\|_2^2)$. This proves the second inequality.
\end{proof}

\section{Geometry of arbitrary maximizing caps near the transition}
Put $w=1+\delta$ and $M_\delta=1+w^2/2$. Let $\overline m(w)$ be the maximum of the cap functional. Its attainment was proved in \texttt{optimality-uniqueness.tex}; every sofa is bounded by this cap maximum. Choose any global maximizing cap $K_\delta$, write its supports as
\[
 p_\delta=p_\delta^*+F_\delta,\qquad k_\delta=k_\delta^*+G_\delta,
\]
and put $C_\delta=|N_w(K_\delta)|-I(\gamma_\delta)$.

\begin{lemma}[Uniform closeness and separated tails]\label{lem:geometry}
For all sufficiently small $\delta$, uniformly over the choice of global maximizing cap,
\[
 0\le M_\delta-\overline m(w)=J_w(F_\delta,G_\delta)+C_\delta=O(\delta^5),
\]
\[
 \|F_\delta'\|_2+\|G_\delta'\|_2=O(\delta^{5/2}).
\]
The supports converge in $C^1$ to the corresponding angle-one relaxed supports as $\delta\downarrow0$. The positive-threshold radial-core and disjoint-tail hypotheses of \texttt{endpoint-tail-functional.tex} therefore hold, and
\[
 C_\delta\ge\mathcal R_w(p_\delta)+\mathcal R_w(\widehat k_\delta),\qquad
 \widehat k_\delta(t)=k_\delta(w-t).
\]
\end{lemma}
\begin{proof}
The feasible recovery family in \texttt{transition-recovery.tex}, or its unperturbed special case, gives $\overline m(w)\ge M_\delta-O(\delta^5)$. The all-angle maximizing-cap majorization gives $C_\delta\ge0$, and the exact relaxation gap gives $J_w\ge0$. Their sum is the displayed cap deficit, hence each is $O(\delta^5)$. Lemma~\ref{lem:local} gives the derivative estimates. The pinned traces imply the same order of uniform support estimates by Cauchy--Schwarz.

The fixed-angle curvature theorem bounds the two active curvature densities uniformly by $\kappa(D)$, where $D$ is a diameter bound for the endpoint parallelograms for $w$ near one. The supports themselves are uniformly bounded. Thus the interior-sided first derivatives are uniformly Lipschitz on their closed active parameter intervals, including the endpoints whose atoms were excluded. Their $L^2$ differences tend to zero, so their uniform differences tend to zero as well: otherwise the uniform Lipschitz bound would produce an interval with a nonvanishing $L^2$ difference. This gives $C^1$ convergence, with the parameter endpoints interpreted on their varying intervals.

At the limiting cap the thresholds are positive on the open interval. Near the two pinned endpoints positivity follows from its nonzero inward derivatives $k_*'(0)=1-c_1$ and $p_*'(1)=c_1-1$, with $c_1=\sec1-\tan1$; the other endpoints have threshold $c_1>0$. Hence the new thresholds stay positive everywhere in the interior. The arms already satisfy the strict all-angle lower bounds. Their corner is therefore a positive radial graph as in the endpoint-tail lemma.

At the limit the niche quadrilaterals lie in a disk of radius $\sqrt2c_1<1/2$. Their vertices depend continuously on the supports and angle, with both cosine denominators bounded away from zero, so a common radius below one still bounds the niches. Simultaneous membership in the two limiting tail half-planes requires ordinate at least one, since $c_1(1+\sin1)/\cos1=1$. The strict gap from the niche disk persists under perturbation. Thus the two tails are disjoint, and their area bound follows from the exact core/tail lemma.
\end{proof}

\section{Blowing up a maximizing cap, not just a trial family}
Given any sequence $\delta\downarrow0$, define functions on the positive half-line, extended by zero past $w/\delta$, by
\[
 v_\delta^R(s)=\delta^{-2}F_\delta'(\delta s),\qquad
 v_\delta^L(s)=-\delta^{-2}G_\delta'(w-\delta s).
\]
Their $L^2$ norms are uniformly bounded by Lemma~\ref{lem:geometry}. Pass to a subsequence with weak limits $v^R,v^L\in L^2(0,\infty)$, and put $q^R(s)=\int_0^s v^R$, $q^L(s)=\int_0^s v^L$.

\begin{lemma}[Lower bound by two exact transition models]\label{lem:liminf}
Along this subsequence,
\[
 \liminf_{\delta\downarrow0}\frac{M_\delta-\overline m(1+\delta)}{\delta^5}
 \ge\mathscr E(v^R)+\mathscr E(v^L),
\]
where $\mathscr E$ is the functional of \texttt{transition-model.tex}.
\end{lemma}
\begin{proof}
The primitives of $v_\delta^R$ converge uniformly on each compact interval to $q^R$. Pointwise convergence follows by testing weak $L^2$ convergence against the indicator of $[0,s]$; uniformity follows from the common square-root modulus of continuity. These primitives equal
\[
 q_\delta^R(s)=\delta^{-3}\bigl(F_\delta(\delta s)-F_\delta(0)\bigr).
\]
The analogous reflected statement holds at the other endpoint.

For the right-tail intercept, uniform support control gives $|F_\delta(0)|=O(\delta^{5/2})$. Therefore, for each fixed $s>0$,
\begin{align*}
 &\delta^{-3}\left(\frac{p_\delta(\delta s)-1}{\cos(\delta s)}-p_\delta(0)+1\right)\\
 &\qquad\longrightarrow q^R(s)+s^2/2-s^3/6.
\end{align*}
Indeed the relaxed part follows from its exact intercept identity, the difference $F_\delta(\delta s)-F_\delta(0)$ is the primitive above, and multiplication by $\sec(\delta s)-1=O(\delta^2)$ adds only $O(\delta^{9/2})$, negligible after division by $\delta^3$.

In the right-tail integral use $z=\delta^3X$. For every fixed test parameter $s$, the tent height divided by $\delta^2$ tends to $(q^R(s)+s^2/2-s^3/6-X)/s$. Taking a countable dense set of test parameters, and then Fatou's lemma for the nonnegative scaled envelope, gives
\[
 \liminf\frac{\mathcal R_w(p_\delta)}{\delta^5}
 \ge\int_0^\infty\left(\sup_{s>0}\frac{q^R(s)+s^2/2-s^3/6-X}{s}\right)_+dX.
\]
No tightness or interchange of an uncontrolled supremum is needed for this lower bound. Reflection supplies the corresponding left-tail inequality.

For the quadratic gap, apply Lemma~\ref{lem:local} with $\eta=R\delta$ for a fixed finite $R$. Its error divided by $\delta^5$ is $O(R\delta)$, by the derivative-energy bound. Weak lower semicontinuity therefore gives
\[
 \liminf\frac{J_w(F_\delta,G_\delta)}{\delta^5}
 \ge\frac12\int_0^R\bigl((v^R)^2+(v^L)^2\bigr).
\]
Let $R\to\infty$. Combining the nonnegative core/tail bound from Lemma~\ref{lem:geometry} with the three lower estimates proves the lemma. This localization is essential: global mixed terms cannot simply be dropped from the quadratic gap.
\end{proof}

\section{The unrestricted value and asymptotic rigidity}
\begin{theorem}[Fifth-order transition deficit]\label{thm:value}
For the unrestricted prescribed-net-angle sofa problem,
\[
 \boxed{m(1+\delta)=1+\frac{(1+\delta)^2}{2}-\frac{229}{1920}\delta^5+o(\delta^5)}
 \qquad(\delta\downarrow0).
\]
The cap maximum $\overline m(1+\delta)$ has the same expansion.
\end{theorem}
\begin{proof}
For any sequence of positive $\delta$ tending to zero, choose maximizing caps and extract the subsequence used above. The exact model certificate gives
\[
 \mathscr E(v^R)+\mathscr E(v^L)\ge2\cdot\frac{229}{3840}=\frac{229}{1920}.
\]
Lemma~\ref{lem:liminf} therefore bounds the cap deficit from below. The actual feasible recovery sofas give the reverse upper bound on the unrestricted deficit, with that same coefficient. Since
\[
 |S_\delta|\le m(1+\delta)\le\overline m(1+\delta)\le M_\delta,
\]
the two deficits are squeezed to the common coefficient. The sequence was arbitrary, proving the two expansions.
\end{proof}

\begin{theorem}[Unique endpoint support-difference profile]
For any choices of global maximizing caps as $\delta\downarrow0$, the primitives
\[
 \delta^{-3}\bigl(F_\delta(\delta s)-F_\delta(0)\bigr),\qquad
 \delta^{-3}\bigl(G_\delta(w-\delta s)-G_\delta(w)\bigr)
\]
converge uniformly on every compact $s$ interval to
\[
 q_*(s)=\begin{cases}
 -3s/16,&0\le s\le3/2,\\
 s^3/12-s^2/4,&3/2\le s\le2,\\
 -1/3,&s\ge2.
 \end{cases}
\]
Their rescaled derivatives converge weakly in $L^2(0,\infty)$ to the model minimizer $v_*$.
\end{theorem}
\begin{proof}
Every weakly convergent subsequence has, by the proved value limit and Lemma~\ref{lem:liminf}, $\mathscr E(v^R)+\mathscr E(v^L)$ at most twice the model minimum. Each term is at least that minimum, so both attain it. The exact model rigidity inequality forces $v^R=v^L=v_*$. Thus all weak subsequential limits agree. Compactness of the primitives on each bounded interval then gives convergence of the whole family, not just of a selected subsequence.
\end{proof}

\begin{remark}[The exact scope of the rigidity claim]
The profile theorem concerns global maximizing caps and their endpoint-relative support differences. It does not assert exact uniqueness of a cap at a fixed angle greater than one. It also does not assert that $F_\delta(0)/\delta^3$ has a determined limit: an additive support displacement varying over the bulk may have a lower-order energy cost, and the leading local model does not fix it. Nor may the all-maximizer geometric hypotheses automatically be transferred to arbitrary near-maximizing caps or sofas.
\end{remark}

\paragraph{Dependencies and review boundary.}
The proof uses the previously written fixed-angle curvature and all-maximizer majorization, cap attainment, the exact quadratic gap, the positive radial-core/tail lemma, the exact transition-model certificate, and the independently feasible recovery sequence. It does not use a numerical optimizer or assume the larger-angle candidate. The full optimizer classification on $1<\omega<\pi/2$ remains open in this workstream. No CI, Lean compilation, or TeX compilation was run.
\end{document}
